\documentclass[../../main.tex]{subfiles}
% \documentclass[a4paper,10pt]{ltjsarticle}
\begin{document}
\rhead{京都大学後期$2006$年$4$番}
\setcounter{equation}{0}
\begin{oframed}
  \documentclass[../../main.tex]{subfiles}
\begin{document}

平面上の点$\mathrm{O}$を中心とし，半径$1$の円周上に相異なる$3$点$\mathrm{A}$，$\mathrm{B}$，$\mathrm{C}$がある．
$\triangle{\mathrm{ABC}}$の内接円の半径$r$は$\dfrac{1}{2}$以下であることをしめせ．

\end{document}

\end{oframed}
\setlength{\columnseprule}{0.4pt}
\begin{multicols}{2}
  {\bf [解答]}
  表記の簡潔さのため，
  $\angle A = \alpha$, $\angle B = \beta$, $\angle C = \gamma$ とおく．
  ただし
  \begin{align}
    \begin{dcases}
      \alpha, \beta, \gamma > 0 \\
      \alpha + \beta + \gamma = \pi
    \end{dcases}\label{eq:1}
  \end{align}
  である．

  \begin{figure}[H]
    \centering
    \begin{tikzpicture}[x=2.5cm, y=2.5cm, >=stealth, every node/.style={font=\small}]
      % 外接円の半径 R=1．A,B,C は中心角 115°,205°,335° の点（一般の三角形の例）
      % 内心 I と内接円の半径 r はこの配置から計算した値 (I=(aA+bB+cC)/(a+b+c))
      \coordinate (O) at (0, 0);
      \coordinate (A) at (115:1);
      \coordinate (B) at (205:1);
      \coordinate (C) at (335:1);
      \coordinate (I) at (-0.2326, 0.0491);
      \def\rin{0.4717}
      \coordinate (Hc) at ($(A)!(I)!(B)$);
      \coordinate (Ha) at ($(B)!(I)!(C)$);
      \coordinate (Hb) at ($(C)!(I)!(A)$);

      % 外接円と内接円
      \draw[thick, blue!85!black] (O) circle (1);
      \draw[thick, teal!85!black] (I) circle (\rin);

      % 三角形 ABC
      \draw[line width=1.2pt] (A) -- (B) -- (C) -- cycle;

      % 内心から各辺への垂線（長さ r）
      \draw[densely dashed, thick, teal!85!black] (I) -- (Hc);
      \draw[densely dashed, thick, teal!85!black] (I) -- (Ha);
      \draw[densely dashed, thick, teal!85!black] (I) -- (Hb);
      \node[teal!85!black, right=1pt, font=\footnotesize] at ($(I)!0.5!(Ha)$) {$r$};

      % 角
      \pic[draw, "$\alpha$", angle eccentricity=1.6, angle radius=0.45cm] {angle = B--A--C};
      \pic[draw, "$\beta$", angle eccentricity=1.6, angle radius=0.45cm] {angle = C--B--A};
      \pic[draw, "$\gamma$", angle eccentricity=1.6, angle radius=0.45cm] {angle = A--C--B};

      % 辺の長さ
      \node[below=1pt] at ($(B)!0.6!(C)$) {$a$};
      \node[above right=0pt] at ($(C)!0.5!(A)$) {$b$};
      \node[left=1pt] at ($(A)!0.5!(B)$) {$c$};

      % 点とラベル
      \node[circle, fill=blue!85!black, inner sep=1.2pt] at (O) {};
      \node[below right=0pt, blue!85!black] at (O) {$O$};
      \node[circle, fill=teal!85!black, inner sep=1.2pt] at (I) {};
      \node[below left=0pt, teal!85!black] at (I) {$I$};
      \node[above left=1pt] at (A) {$A$};
      \node[left=1pt] at (B) {$B$};
      \node[right=1pt] at (C) {$C$};
    \end{tikzpicture}
    \caption{$\triangle ABC$ と外接円（中心 $O$，半径 $R$）・内接円（中心 $I$，半径 $r$）}
    \label{2006-4:fig:1}
  \end{figure}

  $\triangle ABC$に正弦定理を用いて，外接円の半径$R=1$として
  \begin{align}
    \dfrac{a}{\sin \alpha} = \dfrac{b}{\sin \beta} = \dfrac{c}{\sin \gamma} = 2R = 2 \label{eq:2}
  \end{align}
  となる．

  以下$\triangle ABC$ の面積$S$を2通りで表す．
  まず\cref{eq:2}より
  \begin{align}
    S
    & = \dfrac{1}{2}|AB||AC|\sin\alpha                       \\
    & = \dfrac{1}{2} bc \sin\alpha                           \\
    & = \dfrac{1}{2} (2\sin\beta)(2\sin\gamma) \sin\alpha    \\
    & = 2 \sin\alpha\sin\beta\sin\gamma \label{eq:3}
  \end{align}
  である．

  次に，内接円の半径 $r$ ，中心$I$とおくと
  \begin{align}
    S
    & = \triangle IAB + \triangle IBC + \triangle ICA \\
    & =\dfrac{1}{2} r \left(2R \sin \alpha + 2R \sin \beta + 2R \sin \gamma\right) \\
    & =r \left(\sin \alpha + \sin \beta + \sin \gamma\right) \label{eq:4}
  \end{align}
  である．

  以上\cref{eq:3,eq:4}から
  \begin{align}
    r = \dfrac{2 \sin\alpha\sin\beta\sin\gamma}{\sin \alpha + \sin \beta + \sin \gamma} \label{eq:5}
  \end{align}
  だから，この右辺の最大値が$1/2$以下であることを示せば良い．
  \cref{eq:1}を用いて$\gamma$を削除して，
  \begin{align}
    f(x, y) = \dfrac{2\sin x \sin y \sin(x+y)}{\sin x + \sin y + \sin(x+y)}
  \end{align}
  とおく．値域は
  \begin{align}
    0 < x, y, x+y < \pi \label{eq:6}
  \end{align}
  である．倍角公式，および和積公式から
  \begin{align}
    f(x, y)
    & = 2\dfrac{2 \sin \dfrac{x+y}{2} \cos \dfrac{x+y}{2} \cdot \sin x \sin y}{2 \sin \dfrac{x+y}{2} \cos \dfrac{x-y}{2} + 2 \sin \dfrac{x+y}{2} \cos \dfrac{x+y}{2}} \\
    & = 2\dfrac{\cos \dfrac{x+y}{2} \cdot \sin x \sin y}{\cos \dfrac{x-y}{2} + \cos \dfrac{x+y}{2}}
  \end{align}
  となる．さらに倍角公式，和積公式を用いて
  \begin{align}
    f(x,y)
    & = 2\dfrac{\cos \dfrac{x+y}{2} \cdot 2\sin \dfrac{x}{2}\cos\dfrac{x}{2} \cdot 2\sin \dfrac{y}{2}\cos\dfrac{y}{2}}{2 \cos \dfrac{x}{2} \cos \dfrac{y}{2}} \\
    & = 4\cos \dfrac{x+y}{2} \sin \dfrac{x}{2} \sin \dfrac{y}{2}
  \end{align}
  さらに積和公式を用いると
  \begin{align}
    f(x,y)
    & = 2 \sin \dfrac{x}{2} \left[\sin\left(y+\dfrac{x}{2}\right)-\sin\dfrac{x}{2}\right]
  \end{align}
  と変形できる．
  \cref{eq:6}から$\sin \dfrac{x}{2} > 0$だから，
  $x$を固定して$y$を動かした時，
  $f(x, y)$ は
  \begin{align}
    & \sin\left(y+\dfrac{x}{2}\right) = 1   \\
    \therefore
    & y = \dfrac{1}{2}\left(\pi - x \right)
  \end{align}
  で最大値を取る．（\cref{eq:6}を満たしていることに注意．）
  その最大値$f(x)$とおくと，
  \begin{align}
    & f(x) \equiv \max f(x, y) = 2\sin \dfrac{x}{2} \left[1-\sin\dfrac{x}{2}\right] & (0 < x < \pi)
  \end{align}
  である．
  次に$x$を動かすと，
  \begin{align}
    & \sin \dfrac{x}{2} = \dfrac{1}{2} \\
    \therefore
    & x = \dfrac{\pi}{3}
  \end{align}
  で最大値
  \begin{align}
    \max f(x) = f\left(\dfrac{\pi}{3}\right) = \dfrac{1}{2}
  \end{align}
  をとる．
  従って\cref{eq:5}から
  \begin{align}
    r \le \dfrac{1}{2}
  \end{align}
  である．等号成立条件は
  \begin{align}
    \alpha = \beta = \gamma = \dfrac{\pi}{3}
  \end{align}
  である．以上から題意は示された．$\cdots$(答)

  \vspace{10pt}
  {\bf [解説]}
  解答は，もっとも技巧がいらない代わりに計算量が多い方法で解いた．
  この問題はさまざまな解き方があるのでここでいくつか解説して行く．

  \begin{figure}[H]
    \centering
    \begin{tikzpicture}[x=2.5cm, y=2.5cm, >=stealth, every node/.style={font=\small}]
      % 図1と同じ三角形（R=1）．I は内心，r=0.4717，d=|OI|=0.2377
      % P,Q は直線 OI と外接円の交点（P は I 側），M は AI の延長と外接円の交点（弧 BC の中点），
      % D は M の対蹠点，N は I から AB に下ろした垂線の足
      \coordinate (O) at (0, 0);
      \coordinate (A) at (115:1);
      \coordinate (B) at (205:1);
      \coordinate (C) at (335:1);
      \coordinate (I) at (-0.2326, 0.0491);
      \def\rin{0.4717}
      \coordinate (P) at (-0.9784, 0.2067);
      \coordinate (Q) at (0.9784, -0.2067);
      \coordinate (M) at (270:1);
      \coordinate (D) at (90:1);
      \coordinate (N) at ($(A)!(I)!(B)$);

      \draw[thick, blue!85!black] (O) circle (1);
      \draw[thick, teal!85!black] (I) circle (\rin);
      \draw[line width=1.2pt] (A) -- (B) -- (C) -- cycle;

      % 直径 PQ（O, I を通る）と直径 DM
      \draw[densely dashed, thick, gray!70!black] (P) -- (Q);
      \draw[densely dashed, thick, gray!70!black] (D) -- (M);
      % 角の二等分線 AI とその延長 IM
      \draw[thick, orange!85!black] (A) -- (M);
      % 相似な直角三角形 ANI と DBM
      \draw[thick, red!85!black] (I) -- (N);
      \draw[densely dashed, thick, red!85!black] (D) -- (B) -- (M);
      \draw[densely dashed, thick, gray!70!black] (B) -- (I);
      \pic[draw, red!85!black, angle radius=0.2cm] {right angle = A--N--I};
      \pic[draw, red!85!black, angle radius=0.2cm] {right angle = D--B--M};

      % 点とラベル
      \foreach \p/\pos in {A/above left, B/left, C/right, P/left, Q/right, M/below, D/above, N/above left} {
        \node[circle, fill, inner sep=1.1pt] at (\p) {};
        \node[\pos=1pt] at (\p) {$\p$};
      }
      \node[circle, fill=blue!85!black, inner sep=1.1pt] at (O) {};
      \node[below right=0pt, blue!85!black] at (O) {$O$};
      \node[circle, fill=teal!85!black, inner sep=1.1pt] at (I) {};
      \node[above right=0pt, teal!85!black] at (I) {$I$};
    \end{tikzpicture}
    \caption{オイラーの定理の証明に用いる点（$P,Q$ は直線 $OI$ と外接円の交点，$M$ は $AI$ の延長と外接円の交点，$DM$ は直径，$N$ は $I$ から $AB$ への垂線の足）}
    \label{2006-4:fig:2}
  \end{figure}

  おそらく一番簡潔なのは，初等幾何におけるオイラーの定理（またはオイラーの不等式）を用いると一発で求まる．
  オイラーの定理とは，三角形の内接円の半径を $r$，外接円の半径を $R$ とおくとき，
  外心 $O$ と内心 $I$ との距離 $d$ は以下の式で表される．
  \begin{align}
    d^2 = R(R-2r)
  \end{align}
  従って，左辺が$0$以上なので$R-2r\ge 0 \iff r \le R/2=1/2$である．

  % https://manabitimes.jp/math/580
  ここでこのオイラーの定理の証明を与えておこう．
  上図のように各点をおくと，まず外心と内心の定義から
  \begin{align}
    PI\cdot IQ = (|PO|-|OI|)\cdot (|QO|+|OI|) = (R-d)(R+d) \label{eq:7}
  \end{align}
  である．さらに方べきの定理を$PQ,AM$に用いて
  \begin{align}
    PI \cdot IQ = AI \cdot IM \label{eq:8}
  \end{align}
  である．

  % 左に方べき，右にtriangle IBMの二等辺三角形の図
  \begin{figure}[H]
    \centering
    \begin{tikzpicture}[x=1.65cm, y=1.65cm, >=stealth, every node/.style={font=\footnotesize}]
      % 図2と同じ配置（R=1）
      \coordinate (O) at (0, 0);
      \coordinate (A) at (115:1);
      \coordinate (B) at (205:1);
      \coordinate (C) at (335:1);
      \coordinate (I) at (-0.2326, 0.0491);
      \coordinate (P) at (-0.9784, 0.2067);
      \coordinate (Q) at (0.9784, -0.2067);
      \coordinate (M) at (270:1);

      % 左：方べきの定理（弦 PQ と AM が I で交わる）
      \begin{scope}
        \draw[thick, blue!85!black] (O) circle (1);
        \draw[gray!70!black] (A) -- (B) -- (C) -- cycle;
        \draw[line width=1.3pt, teal!85!black] (P) -- (Q);
        \draw[line width=1.3pt, orange!85!black] (A) -- (M);
        \foreach \p/\pos in {A/above left, P/left, Q/right, M/below, B/left, C/right} {
          \node[circle, fill, inner sep=1pt] at (\p) {};
          \node[\pos=1pt] at (\p) {$\p$};
        }
        \node[circle, fill=red!85!black, inner sep=1.2pt] at (I) {};
        \node[above right=0pt, red!85!black] at (I) {$I$};
        \node[below=3pt, font=\small] at (0, -1.15) {$PI\cdot IQ=AI\cdot IM$};
      \end{scope}

      % 右：二等辺三角形 IBM
      \begin{scope}
        \coordinate (O2) at ($(O)+(2.7,0)$);
        \coordinate (A2) at ($(A)+(2.7,0)$);
        \coordinate (B2) at ($(B)+(2.7,0)$);
        \coordinate (C2) at ($(C)+(2.7,0)$);
        \coordinate (I2) at ($(I)+(2.7,0)$);
        \coordinate (M2) at ($(M)+(2.7,0)$);
        \draw[thick, blue!85!black] (O2) circle (1);
        \draw[gray!70!black] (A2) -- (B2) -- (C2) -- cycle;
        \draw[gray!70!black] (A2) -- (M2);
        \fill[pattern=north east lines, pattern color=red!70!black] (I2) -- (B2) -- (M2) -- cycle;
        \draw[line width=1.3pt, red!85!black] (I2) -- (B2) -- (M2) -- cycle;
        \pic[draw, red!85!black, angle radius=0.3cm] {angle = B2--I2--M2};
        \pic[draw, red!85!black, angle radius=0.3cm] {angle = M2--B2--I2};
        \node[circle, fill, inner sep=1pt] at (A2) {}; \node[above left=1pt] at (A2) {$A$};
        \node[circle, fill, inner sep=1pt] at (B2) {}; \node[left=1pt] at (B2) {$B$};
        \node[circle, fill, inner sep=1pt] at (C2) {}; \node[right=1pt] at (C2) {$C$};
        \node[circle, fill, inner sep=1pt] at (M2) {}; \node[below=1pt] at (M2) {$M$};
        \node[circle, fill=red!85!black, inner sep=1.2pt] at (I2) {}; \node[above right=0pt, red!85!black] at (I2) {$I$};
        \node[below=3pt, font=\small] at ($(O2)+(0,-1.15)$) {$\angle BIM=\angle IBM$};
      \end{scope}
    \end{tikzpicture}
    \caption{左：弦 $PQ$, $AM$ に対する方べきの定理．右：$\angle BIM=\angle IBM$ より $\triangle IBM$ は $BM=IM$ の二等辺三角形}
    \label{2006-4:fig:3}
  \end{figure}

  次に，三角形$IBM$について
  \begin{align}
    \angle BIM
    = \angle BAI + \angle ABI
    = \angle MAC + \angle IBC
    = \angle IBM
  \end{align}
  だから$BM=IM$なる二等辺三角形である．よって
  \begin{align}
    AI \cdot IM = AI \cdot BM \label{eq:9}
  \end{align}
  である．

  最後に$\triangle ANI,\triangle DBM$に注目すると，円周角の定理より
  \begin{align}
    \angle NAI = \angle BDM
  \end{align}
  である．内心の性質より
  \begin{align}
    \angle ANI = \dfrac{\pi}{2}
  \end{align}
  であり，円周角の定理より
  \begin{align}
    \angle DBM = \dfrac{\pi}{2}
  \end{align}
  だから
  \begin{align}
    \angle ANI = \angle DBM
  \end{align}
  である．以上から二角相等より
  \begin{align}
    \triangle ANI \sim \triangle DBM
  \end{align}
  であって相似から
  \begin{align}
    AI \cdot BM = DM \cdot NI \label{eq:10}
  \end{align}
  である．\cref{eq:7,eq:8,eq:9,eq:10}をまとめて
  \begin{align}
    (R-d)(R+d)
    = PI \cdot IQ
    = AI \cdot IM \\
    = AI \cdot BM
    = DM \cdot NI = 2Rr
  \end{align}
  より示された．

  次に，イェンゼンの不等式を用いた方法を紹介する．
  \cref{eq:5}以降，回答では文字消去をして解いていったが，
  当然対称なまま変形していった方が綺麗に変形していける．
  まず，倍角公式より分子は
  \begin{align}
    \text{分子}
    & =\sin\alpha\sin\beta\sin\gamma                                                                                            \\
    & = 8\sin\dfrac{\alpha}{2}\cos\dfrac{\alpha}{2}\sin\dfrac{\beta}{2}\cos\dfrac{\beta}{2}\sin\dfrac{\gamma}{2}\cos\dfrac{\gamma}{2}
  \end{align}
  と変形できる．
  次に分母は和積公式を繰り返し用いて変形する．
  途中で$\alpha+\beta+\gamma=\pi$を利用する．
  \begin{align}
    \text{分母}
    & = \sin \alpha + \sin \beta + \sin \gamma                                                            \\
    & = 2\sin\dfrac{\alpha+\beta}{2}\cos\dfrac{\alpha-\beta}{2} + 2\sin\dfrac{\gamma}{2}\cos\dfrac{\gamma}{2} \\
    & = 2\sin\dfrac{\pi-\gamma}{2}\cos\dfrac{\alpha-\beta}{2} + 2\sin\dfrac{\gamma}{2}\cos\dfrac{\gamma}{2}   \\
    & = 2\cos\dfrac{\gamma}{2}\cos\dfrac{\alpha-\beta}{2} + 2\sin\dfrac{\gamma}{2}\cos\dfrac{\gamma}{2}       \\
    & = 2\cos\dfrac{\gamma}{2}\left[\cos\dfrac{\alpha-\beta}{2} + \sin\dfrac{\gamma}{2}\right]               \\
    & = 2\cos\dfrac{\gamma}{2}\left[\cos\dfrac{\alpha-\beta}{2} + \sin\dfrac{\pi-\alpha-\beta}{2}\right]     \\
    & = 2\cos\dfrac{\gamma}{2}\left[\cos\dfrac{\alpha-\beta}{2} + \cos\dfrac{\alpha+\beta}{2}\right]         \\
    & = 2\cos\dfrac{\gamma}{2}\cdot 2\cos\dfrac{\alpha}{2}\cos\dfrac{\beta}{2}                               \\
    & = 4\cos\dfrac{\alpha}{2}\cos\dfrac{\beta}{2}\cos\dfrac{\gamma}{2}
  \end{align}
  となる．従って，
  \begin{align}
    r = 4R\sin\dfrac{\alpha}{2}\sin\dfrac{\beta}{2}\sin\dfrac{\gamma}{2}
  \end{align}
  となる．各項は正だから，相加相乗平均より
  \begin{align}
    r \le 4R \left(\dfrac{\sin\dfrac{\alpha}{2}+\sin\dfrac{\beta}{2}+\sin\dfrac{\gamma}{2}}{3}\right)^3
  \end{align}
  さらに$\sin x$は$0<x<\pi$で上に凸だから，イェンゼンの不等式より
  \begin{align}
    r
    & \le 4R \left[\sin\left(\dfrac{\dfrac{\alpha}{2}+\dfrac{\beta}{2}+\dfrac{\gamma}{2}}{3}\right)\right]^3 \\
    & = 4R\left[\sin\dfrac{\pi}{6} \right]^3                                                              \\
    & = \dfrac{R}{2}                                                                                      \\
    & = \dfrac{1}{2}
  \end{align}
  となる．よって示された．
  等号成立は$\alpha=\beta=\gamma=\pi/3$のときである．

  以上の例は，角度を用いて変形していく方法だったが，
  辺の長さ$a,b,c$を用いて変形して行く方法もある．
  \cref{eq:3,eq:4}を長さを用いて書くと
  \begin{align}
    S = \dfrac{abc}{4R} = \dfrac{1}{2}r(a+b+c)
  \end{align}
  である．さらに，ヘロンの公式から
  \begin{align}
    S = \dfrac{1}{4}\sqrt{(a+b+c)(-a+b+c)(a-b+c)(a+b-c)}
  \end{align}
  とも書けることに留意する．
  従って，
  \begin{align}
    R-2r
    & = \dfrac{abc}{4S} - \dfrac{4S}{a+b+c}    \\
    & = \dfrac{abc(a+b+c) - 16S^2}{4S(a+b+c)}
  \end{align}
  この分子$P$が0以上であることを示す．ヘロンの公式を代入して
  \begin{align}
    P
    & = abc(a+b+c) - 16S^2                             \\
    & = abc(a+b+c) - (a+b+c)(a+b-c)(b+c-a)(c+a-b)      \\
    & = (a+b+c)\left[abc-(a+b-c)(b+c-a)(c+a-b) \right]
  \end{align}
  である．ここで，Ravi変換と呼ばれる変換
  \begin{align}
    x & = \dfrac{a+b-c}{2} \\
    y & = \dfrac{b+c-a}{2} \\
    z & = \dfrac{c+a-b}{2}
  \end{align}
  を代入する．$a,b,c$が三角形の辺だから$x,y,z>0$である．
  \begin{align}
    P
    & = (a+b+c)\left[(x+y)(y+z)(z+x) - 8xyz\right]                         \\
    & = (a+b+c)\left[(xyz+xyz+x^2y+x^2z+y^2x+y^2z+z^2x+z^2y) - 8xyz\right]
  \end{align}
  ここで，大括弧の中は相加相乗平均より
  \begin{align}
    & \dfrac{xyz+xyz+x^2y+x^2z+y^2x+y^2z+z^2x+z^2y}{8} \ge \sqrt[8]{(xyz)^8} = xyz \\
    \therefore
    & (xyz+xyz+x^2y+x^2z+y^2x+y^2z+z^2x+z^2y) - 8xyz \ge 0
  \end{align}
  だから，$P\ge 0 \iff R \ge 2r$である．よって題意は示された．

  \newpage
\end{multicols}
\end{document}
